\section{Compound AI Systems 的范式}

\subsection{什么是 Compound AI Systems}

\begin{knowledgebox}{复合 AI 系统 vs. 单模型}
复合 AI 系统由多个相互作用的组件组成：LLM 调用、检索器、代码执行器、验证器等。每个组件有自己的签名（输入/输出类型），通过程序逻辑组合。这类似于传统软件工程中的``函数组合''。
\end{knowledgebox}

\begin{itemize}
    \item \textbf{RAG} 是最简单的复合系统：检索 + LLM 合成
    \item \textbf{ReAct} 是循环式复合系统：推理 + 工具调用的迭代
    \item \textbf{Multi-hop QA}：多次检索和推理的串联
    \item 更复杂的系统：多阶段 pipeline，包含验证、反思等步骤
\end{itemize}

\subsection{手工 Prompting 的问题}

\begin{warningbox}{Prompting 不可持续}
手工编写 prompt 看似简单，但在复合系统中会导致：
\begin{itemize}
    \item \textbf{脆弱性}：更换模型或修改 pipeline 需要重写所有 prompt
    \item \textbf{不可组合}：难以将多个 prompt 模块化组合
    \item \textbf{难以优化}：无法系统性地搜索最优 prompt
    \item \textbf{不可移植}：一个模型的好 prompt 对另一个模型可能无效
\end{itemize}
\end{warningbox}

\subsection{本章小结}

手工 prompting 在复合系统中不可持续，需要更抽象的编程范式。

